\begin{exercisechunk}
\item \textbf{Compose two input-specified lookup tables.}
\label[exercise]{ex:l14-causal-lookup}
\exerciseprofile{2}{2}{2}{\exproof\quad\excalculation\quad\exconstruction\quad\exinterpretation}{%
Use the output of one associative lookup as the query for a second table.}
Let
\[
 f:\{\mathtt a,\mathtt b\}\to\{1,2\},
 \qquad
 g:\{1,2\}\to\{\mathtt{red},\mathtt{blue}\}
\]
be specified by
\[
 f(\mathtt a)=2,\quad f(\mathtt b)=1,
 \qquad
 g(1)=\mathtt{red},\quad g(2)=\mathtt{blue}.
\]
The input sequence contains the tagged records
\[
 ((\mathtt F,\mathtt a,2),(\mathtt F,\mathtt b,1),
  (\mathtt G,1,\mathtt{red}),(\mathtt G,2,\mathtt{blue}),
  (\mathtt{QUERY},\mathtt b)).
\]
In one lookup round, the final position may form a query from its current
state, equality-select one earlier tagged record, retrieve that record's
value, and update its state. All lookups must be causal.
\begin{enumerate}[(a),beginpenalty=10000]
\item List the query, selected record, retrieved value, and updated final
state in each of two rounds that return $g(f(\mathtt b))$. Compute the final
answer.
\item Let $\cD,\cV,\cW$ be finite sets, let $f:\cD\to\cV$ and
$g:\cV\to\cW$, and suppose the input contains tagged records for both maps
followed by a query $x\in\cD$. Assume the keys in each table are distinct.
Give one two-round lookup program and prove that its output is $g(f(x))$ for
every choice of the two tables and every allowed set size.
\item Explain why the final query position cannot in general return
$g(f(x))$ in one round under the stated one-record-per-round abstraction.
Identify one change to the input representation that would make one-round
retrieval possible.
\end{enumerate}
\end{exercisechunk}
